% !TEX root = ../main.tex

% =========================================================
% 致谢
% 规范要求：
% 1. 装订顺序：置于附录之后、封底之前
% 2. 标题格式：同一级标题格式（二号华文中宋加粗，居中，段前3行，段后2行）
% 3. 内容格式：同正文格式（小四号宋体，1.5倍行距，首行缩进2字符）
% 4. 篇幅限制：内容严格限一页
% 5. 盲审控制：盲审模式下自动匿名或隐藏致谢内容
% =========================================================

\ifswufereview
  % 盲审模式下自动跳过致谢页，避免泄露导师或个人信息
\else
  \phantomsection
  \addcontentsline{toc}{chapter}{致谢}
  \chapter*{致\quad 谢}
  \markboth{致谢}{致谢}

  时光荏苒，白驹过隙。转眼间在西南财经大学的四年本科求学生涯即将画上圆满的句号。回首在柳林校区的求学岁月，从草堂明月到腾飞广场，从光华铁树到晨曦微露的图书馆，每一次探索与求知都让我受益终生。

  在此，我要特别感谢我的指导教师李四教授。在毕业论文的选题论证、方案设计、理论推导以及论文定稿的全过程中，导师倾注了大量的心血，给予了我悉心细致的指导与由衷的关怀。导师严谨治学的学术风范、求真务实的治学态度和宽厚仁厚的长者风度，将成为我未来人生与学术征程中永远的指路明灯。

  感谢学院各位授课教师在大学四年中传道授业解惑，为我打下了坚实厚重的数理基础与专业素养。感谢同窗好友与课题组各位学长学姐在研究过程中的深入探讨与无私帮助。

  最后，深深感谢我的父母与家人。正是你们二十余年来无微不至的关怀、包容与坚定支持，才使我能够心无旁骛、踏实前行。

  谨以此篇论文，向所有关怀、支持和帮助过我的人致以最崇高的敬意与最诚挚的感谢！
\fi
